\section{Submission Instructions}

\subsection{Methods Paper—maximum word count: 6500}

Methods Papers present an advance likely to make major impact in one or more applied
fields. While theory (e.g. theorems concerning statistical or
learning-theoretic properties) are welcome, this is not essential, but intuition and
understanding of why the method works is important.
We are also very open to theory in a broader sense including examples and conjectures.
Generally, we would expect empirical results to demonstrate effectiveness.
However, in cases where the extent of methodological/conceptual innovation is large
enough, results themselves may be illustrative and need not be of immediate high impact
in any particular applied domain.

Each piece should include:

\begin{itemize}
  \item Unstructured Abstract—maximum word count: 250
  \item Keywords—maximum of 6 and minimum of 5
  \item May include tables and figures—no limit
  \item May include the following back matter sections:
  Acknowledgements, Author contributions (with CRediT details), Conflicts of interest, Funding
  \item Must include the following:
  Data availability, References—no limit
  \item Each submission must contain the following sections and use these terms as the first
  level section headers: Introduction, Methods, Results, Discussion and Conclusion
  (Discussion and Conclusion may be combined).
\end{itemize}


\subsection{Applications Paper—maximum word count: 6500}

Applications Papers are expected to show clear potential for impact in one or more applied domains. Methods should be interesting and appropriate but need not be as novel as for methods papers. Work should be presented so as to be accessible to a broad data science audience, with appropriate explanation of domain-specific ideas. At the same time, clear exposition of the core data scientific ideas involved in the method or application is important. Befitting the focus on data science per se, appropriate mathematical notation would normally be expected to be included in the main text.



Each piece should include:
TBD


\subsection{General requirements}
Below, we provide some general guidance on submission. Our aim is to make submission as straightforward as possible, but request that authors read through the brief points below in preparing papers for submission. See the Manuscript preparation section below for more detailed information on formatting, structure, and style.

Straightforward formatting requirements
We aim to make first submission as straightforward as possible and impose only minimal formatting requirements at this stage. For latex users, you are welcome to use our latex style file, but this is not mandatory: any standard format (e.g. article) and any conventional referencing style is acceptable. Note that while we are happy to accept submissions in any standard format, upon acceptance papers must be reformatted using our style files.

Writing style and use of technical terms
RSS: Data Science and Artificial Intelligence is aimed at a broad readership. However, unlike broad audience journals spanning all scientific fields, the focus is on the data sciences per se. At the same time, in contrast to venues focused on a particular field or subfield (AI, biostatistics, econometrics), results should be presented in such a way as to be understandable by readers from a wide range of (data science) fields. We provide some guidance below on achieving a good balance:

Technical material
RSS: Data Science and Artificial Intelligence is a data science journal and welcomes data scientific submissions with technical content (computational, mathematical, theoretical, conceptual…). While our audience is broad, most readers will be familiar with core, common concepts in data science, hence such notions can be assumed as background. That said, terms and ideas that are more specific to a specific field or subfield should be briefly explained.

Use of field-specific terminology
Terms and ideas that are not standard across the data sciences should be briefly explained on first appearance. Some examples of terms/ideas that are clearly standard background are: estimator, law of large numbers, cross-validation, PCA, posterior density, graphical model. Examples of terms that are more specific (relevant fields or subfields named in parentheses) are: contrastive learning (AI), inverse probability weighting (statistics), skip connection (AI), two-stage least squares (econometrics, statistics), gene regulatory network (bioinformatics), regression discontinuity design (econometrics, statistics). For terms of the latter kind, a brief explanation on first appearance will help make your ideas accessible also to readers from other fields, e.g. “a gene regulatory network is a graph whose nodes are identified with genes and edges indicate biological or biochemical relationships between genes. Here we focus on networks with the following semantics…”

Main text and appendices
We encourage presentation of data scientific material, including technical content, in the main text. We are open to various paper lengths. That said, in the interests of exposition, it may be useful to move some details to Appendices. On first submission, we generally leave these choices (length, main text vs Appendix) to authors, with the only constraint being that the main text should not greatly exceed the length of a conventional journal paper (6500 words).

Figures and tables
In communicating to a relatively broad audience, overview/workflow type figures can be useful. Similarly, relatively detailed figure captions are welcome, since they can help readers understand a key idea in a self-contained fashion and complement details in the text itself. We do not place any hard constraint on the number of figures or tables, other than the general, overall length constraint above. Please see below for further information on figures and further information on tables.


\subsection{Instructions Resources}

\href{https://academic.oup.com/rssdat/pages/general-instructions}{RSS Instructions} 

\href{https://academic.oup.com/pages/authoring/books/preparing-your-manuscript/working-in-latex}{RSS Working in \LaTeX}.
